\section{Introduction}

Deep search agents, such as OpenAI Deep Research and Tongyi DeepResearch \citep{Team2025TongyiDeepResearchTechnical}, are LLM-based systems for long-horizon information seeking and answering complex and multi-hop questions, where evidence is scattered across documents and must be searched for, checked, and synthesized before a conclusion holds \citep{Shi2025DeepResearchA}. The agent reasons, searches, reads what comes back, and decides turn by turn what to look for next, until it can commit to an answer \citep{Team2025MiroThinkerPushingthe}. Over one question, an agent may sustain retrieval and reasoning over tens of steps, and this long-horizon capacity is what separates it from single-shot retrieval-augmented generation.

However, evaluation practice has not caught up with these advances. Progress is still reported primarily through final-answer metrics, such as answer accuracy on BrowseComp \citep{Wei2025BrowseCompASimple}, while the trajectory that produces the answer remains an opaque byproduct. 

This poses concrete challenges for diagnosing agent behavior and developing improvement strategies:
\textbf{1) Wasted effort is mistaken for capability.} Long horizons and heavy tool use are routinely showcased as strengths, yet an effective and efficient search agent should instead adapt its search effort to what the problem genuinely requires: searching more only when the problem calls for it, not as a default behaviour.
\textbf{2) Bottlenecks get misdiagnosed.} A wrong answer may read as ``the agent could not find the evidence,'' so the fix is more or better search. Yet some errors occur with the decisive evidence already in hand, and more search cannot repair them. Differentiating between error types is essential for prescribing the right fix. 
\textbf{3) Persistence and floundering look the same.} Trajectories often run long after new evidence stops arriving, yet the score cannot tell late turns that are closing in on an answer from those that never anchor on the right direction, the very distinction between needing to search \emph{deeper} and needing to search \emph{better}.

To bridge this gap, we propose a diagnostic framework that opens up the trajectory itself. 
By grounding every step in human-annotated, document-level relevance, it measures retrieval-side and reasoning-side behaviour directly, and analyzes search agents from coarse to fine granularity.
Specifically, we organize our analysis into three key perspectives:
\textbf{1) Outcome attribution}: which signal actually tracks final accuracy: search effort or gold-evidence recall?
\textbf{2) Failure-mode taxonomy}: every error is sorted into a \emph{retrieval gap} (supporting evidence never found) or a \emph{utilization gap} (evidence found but still incorrect).
\textbf{3) Trajectory mechanism and query strategy}: each trajectory is followed episode by episode, tracking how the agent searches and when and where evidence arrives.

We conduct our analysis with six agents, spanning two scale tiers: three mid-scale (Tongyi DeepResearch \citep{Team2025TongyiDeepResearchTechnical}, Qwen3.5-35B-A3B \citep{qwen2025qwen3}, gpt-oss-120b \citep{openai2025gptoss}) and three frontier-scale (Kimi K2.6 \citep{kimi2026k26}, GLM 5.1 \citep{zeng2026glm}, Deepseek V4 Pro \citep{deepseekai2026deepseekv4}). All are run under the same ReAct framework, tools, and retriever on BrowseComp-Plus \citep{Chen2025BrowseComp-PlusAMore}, whose fixed corpus and human-annotated qrels supply the relevance grounding; we further validate the main conclusions on BrowseComp \citep{Wei2025BrowseCompASimple} with open-web search.

Our analysis distills into five findings, which trace a single thread from symptom through mechanism to prescription: 

\textbf{1) More effort does not buy more accuracy.} Across the six agents we analyze, neither search volume nor context spend predicts accuracy; if anything, the more accurate agents are also the cheaper ones (§\ref{sec:results-attribution}).

\textbf{2) What does track accuracy is retrieval recall.} Accuracy follows how much of the gold evidence an agent's queries cumulatively retrieve; on the same question, the agent that retrieves the gold evidence answers correctly most of the time, while the agent that misses it almost never does (§\ref{sec:results-attribution}, §\ref{sec:results-synthesis}).

\textbf{3) Failures split into two gaps with opposite fixes.} When agents do fail, the error is either a \emph{retrieval gap} --- the evidence was never surfaced --- or a \emph{utilization gap} --- it was surfaced, yet the answer still went wrong. Agents at nearly identical accuracy can sit at opposite ends of this balance, so the fix is per-diagnosis rather than per-leaderboard: better queries repair the former, better verification the latter (§\ref{sec:results-failure}).

\textbf{4) Evidence early or never, then a wasted tail.} The great majority of search episodes add no new evidence, and the waste is not spread evenly: whatever evidence a trajectory will surface, it surfaces in the early episodes. Correct runs assemble their evidence early and stop soon after; incorrect runs surface only part of it, or none at all, yet keep searching far longer without recovering the rest --- their extra length is a wasted tail, not late progress (§\ref{sec:results-mechanism}).

\textbf{5) Stronger agents search cleaner, not more.} The behavioral signature of strength is query discipline: redundant re-queries are rare but mark failure. Query strategy is not decisive: agents favor different styles, from broad pivoting to anchor-and-verify, and what counts is whether queries \emph{land} on evidence (§\ref{sec:results-query}).

Together, the findings argue that the next gains lie not in deeper search but in better-directed search (§\ref{sec:discussion}):
\textbf{1) For weaker agents, fix the queries:} improve query direction and search strategy, so that the agent recognizes when a hypothesis has failed and pivots to a new angle, instead of re-asking near-paraphrases of queries that already failed.
\textbf{2) For stronger agents, shift effort to verification:} the strongest agent already sits near the oracle-reader ceiling; the headroom that remains lies in verification and answer normalization, not in more retrieval.
\textbf{3) For both, improve at the harness level:} snippet-stream management so old results do not dominate the context window, visit gating against low-value page opens, and evidence-driven stopping in place of fixed budgets.

The paper contributes (i) a diagnostic framework for search agents grounded in relevance; (ii) a controlled comprehensive study of six agents and the five findings above, tracing search behavior and failure modes from coarse to fine; and (iii) a planned release of trajectories and code to support follow-up analysis.
